% ID: SEL-D0025-OA-E01
% ESTADO: cotejada. Segunda pasada visual completa y borrador compilado revisados.
% Fuente: Selectividad/matematicas_aplicadas/suplemente_junio_EXAMEN 4-ANDALUCÍA.pdf
% SHA-256: c1a3a7a37d38b04359115c64762b1f8a543e4d66155338931de84ef6c304da2c
% Página: 1; opción A; ejercicio original 1.

\textbf{(2.5 puntos)} Una empresa envasa y comercializa leche entera y leche desnatada. El litro de leche entera envasado genera un beneficio diario a la empresa de 0.4 \texteuro{} y el de leche desnatada de 0.1 \texteuro{}. La tecnología de la empresa impone que el número de litros de leche entera que se envasan diariamente no supere el doble del número de litros de leche desnatada. Además, la cantidad máxima de leche que se puede envasar diariamente es un total de 3000 litros y solo se dispone de 1200 litros diarios de leche entera para envasar.

\textquestiondown{}Cuánto debe envasar de cada producto para obtener el beneficio máximo? \textquestiondown{}A cuánto ascendería este beneficio?
